\section{Broader Authoritative-State Externalisation Principle}
\label{sec:authoritative-principle}

RaS may be a software-engineering-specific instance of a broader systems idea:

\begin{quote}
\textbf{Persist authoritative state. Reconstruct computation.}
\end{quote}

Suppose a domain has an authoritative state $A_t$, historical interaction $H_t$, a new task $U_t$, and a useful future action $X_{t+1}$. A broader externalisation hypothesis would ask whether:
\begin{equation}
P(X_{t+1}\mid H_t,A_t,U_t)
\approx
P(X_{t+1}\mid A_t,U_t).
\label{eq:authoritative-general}
\end{equation}

If Equation~\ref{eq:authoritative-general} held for a domain, long-term model-native history could have limited marginal value relative to reconstructing from authoritative domain state. Candidate stores might include CRM systems, transactional databases, ticket systems, event logs, document systems, or infrastructure state.

This is future work, not a result of the present paper. Software repositories have several properties that make them unusually favourable: versioning, explicit diffs, executable tests, deterministic build systems, inspectable schemas, mature access control, and a long engineering culture of encoding decisions into artefacts. Other domains may lack equivalent verification or may depend heavily on tacit, interpersonal, or continuously changing state.

There is also no reason to assume that a single authoritative store is sufficient. A production software programme may require repository state plus an issue tracker, deployment state, incident history, and external requirements. The broader principle should therefore be understood as an invitation to identify the \emph{authoritative domain state} rather than to force every agent into Git.

The proposed generalisation becomes scientifically useful only if it creates falsifiable questions: what state must survive, what computation can be safely reconstructed, how much reconstruction costs, and what quality is lost when model-native history is removed? RaS provides one concrete domain in which those questions can be asked experimentally.
